\documentclass[11pt,a4paper]{article}

% ============================================================
%  QIMMA -- COURS : Nombres complexes (2) — forme trigonométrique,
%  exponentielle et applications géométriques
%  2ème Bac Sciences Mathématiques (A et B)
%  Standard exhaustif : définitions formelles + démonstrations
%  Basé sur templates/qimma-template.tex (charte Qimma)
% ============================================================

\usepackage[french]{babel}
\usepackage[utf8]{inputenc}
\usepackage[T1]{fontenc}
\usepackage{lmodern}
\usepackage[a4paper,margin=2cm,top=2.3cm,bottom=2.6cm]{geometry}
\usepackage{amsmath,amssymb}
\usepackage{xcolor}
\usepackage{graphicx}
\usepackage{tikz}
\usetikzlibrary{shadows,calc}
\usepackage[most]{tcolorbox}
\usepackage{titlesec}
\usepackage{fancyhdr}
\usepackage{enumitem}
\usepackage{pifont}
\usepackage{array}
\usepackage{colortbl}
\usepackage{hyperref}

% ------------------- CHARTE QIMMA -------------------
\definecolor{qteal}{HTML}{0d9488}
\definecolor{qtealdark}{HTML}{134e4a}
\definecolor{qamber}{HTML}{fbbf24}
\definecolor{qambertext}{HTML}{92400e}
\definecolor{ink}{HTML}{1f2937}
\definecolor{inkgray}{HTML}{6b7280}
\definecolor{tealpale}{HTML}{e6f5f3}
\colorlet{colDef}{qteal}
\definecolor{colProp}{HTML}{0f766e}
\definecolor{colEx}{HTML}{b45309}
\definecolor{colRem}{HTML}{9d174d}
\color{ink}

\graphicspath{{../../../../assets/}}
\newcommand{\logofile}{../../../../assets/qimma-logo.pdf}

% ------------------- METADONNEES -------------------
\newcommand{\matiere}{Mathématiques}
\newcommand{\niveau}{2\up{ème} Bac -- Sciences Mathématiques}
\newcommand{\chapitretitre}{Nombres complexes (2) -- forme trigonométrique}

% ------------------- EN-TETE / PIED DE PAGE -------------------
\pagestyle{fancy}
\fancyhf{}
\renewcommand{\headrulewidth}{0pt}
\renewcommand{\footrulewidth}{0.5pt}
\fancyfoot[L]{\raisebox{-0.15cm}{\includegraphics[height=0.35cm]{\logofile}}~\small\sffamily\color{inkgray}Qimma}
\fancyfoot[C]{\small\sffamily\color{inkgray}\matiere\ -- \niveau}
\fancyfoot[R]{\small\sffamily\color{qtealdark}\thepage}

% ------------------- TITRES DE SECTION -------------------
\titleformat{\section}
  {\normalfont\sffamily\Large\bfseries\color{qtealdark}}
  {\tikz[baseline=-3.2pt]{\node[circle,fill=qteal,text=white,inner sep=0pt,minimum size=8mm]{\sffamily\bfseries\thesection};}}
  {10pt}{}
  [\vspace{-2mm}{\color{qamber}\titlerule[1.6pt]}\vspace{2mm}]
\titlespacing*{\section}{0pt}{16pt}{10pt}

\titleformat{\subsection}
  {\normalfont\sffamily\large\bfseries\color{qtealdark}}
  {\color{qteal}\ding{212}}{6pt}{}
\titlespacing*{\subsection}{0pt}{12pt}{6pt}

% ------------------- BOITES PEDAGOGIQUES -------------------
\newtcolorbox{fiche}[3][]{
  enhanced, breakable,
  colback=white, colframe=#2,
  boxrule=1pt, arc=2mm,
  top=7mm, left=4mm, right=4mm, bottom=3mm,
  drop shadow={black!25, shadow xshift=1mm, shadow yshift=-1mm},
  attach boxed title to top left={xshift=4mm, yshift=-3mm, yshifttext=-1mm},
  boxed title style={colback=#2, arc=1.5mm, boxrule=0pt, left=2mm, right=2mm},
  coltitle=white, fonttitle=\sffamily\bfseries,
  title={#3}, #1
}
\newenvironment{definition}[1][Définition]
  {\begin{fiche}{colDef}{\ding{110}~~#1}}{\end{fiche}}
\newenvironment{propriete}[1][Propriété]
  {\begin{fiche}{colProp}{\ding{72}~~#1}}{\end{fiche}}
\newenvironment{exemple}[1][Exemple]
  {\begin{fiche}{colEx}{\ding{217}~~#1}}{\end{fiche}}
\newenvironment{remarque}[1][Remarque]
  {\begin{fiche}{colRem}{\ding{43}~~#1}}{\end{fiche}}

% Démonstration (hors boîte, style discret)
\newenvironment{demo}
  {\par\smallskip\noindent{\sffamily\bfseries\color{qtealdark}Démonstration.}\ \itshape}
  {\ \normalfont\hfill$\blacksquare$\par\medskip}

\hypersetup{colorlinks=true, linkcolor=qtealdark, urlcolor=qtealdark}

% ============================================================
\begin{document}

% ------------------- BANNIERE DE TITRE -------------------
\begin{tcolorbox}[
  enhanced, boxrule=0pt, arc=4mm,
  interior style={left color=qtealdark, right color=qteal},
  top=8mm, bottom=8mm, left=10mm, right=32mm,
  drop shadow={black!35, shadow xshift=1.5mm, shadow yshift=-1.5mm},
  before skip=0pt, after skip=9mm,
  overlay={
    \node[anchor=north west] at ([xshift=6mm,yshift=-6mm]frame.north west)
      {\includegraphics[height=1.25cm]{\logofile}};
    \node[anchor=north east, fill=qamber, text=qambertext, rounded corners=2pt,
          inner sep=4pt, font=\sffamily\bfseries\small] at ([xshift=-6mm,yshift=-6mm]frame.north east) {COURS};
  }
]
\hspace{1.6cm}\centering
{\sffamily\bfseries\color{white}\fontsize{23}{27}\selectfont \chapitretitre}\\[3mm]
{\sffamily\color{white!90}\large \matiere\ \textbullet\ \niveau}
\end{tcolorbox}

% ------------------- OBJECTIFS -------------------
\begin{tcolorbox}[
  enhanced, colback=tealpale, colframe=qteal, boxrule=0.8pt, arc=2mm,
  left=5mm, right=5mm, top=3mm, bottom=3mm,
  title={\sffamily\bfseries\color{qtealdark}\ding{51}~~Objectifs du chapitre},
  coltitle=qtealdark, colbacktitle=tealpale, boxrule=0.8pt,
  attach title to upper={\par\vspace{1mm}}
]
\begin{itemize}[leftmargin=6mm, label=\ding{212}, itemsep=1mm]
  \item Déterminer un argument, écrire un complexe sous forme trigonométrique et exponentielle, et passer d'une forme à l'autre.
  \item Démontrer et utiliser les propriétés de l'argument (produit, quotient, puissance) et la formule de Moivre.
  \item Utiliser les formules d'Euler pour linéariser ou développer des expressions trigonométriques.
  \item Calculer les racines $n$-ièmes d'un nombre complexe, en particulier les racines de l'unité.
  \item Exploiter les complexes en géométrie : angle orienté, alignement, orthogonalité, cocyclicité de quatre points, similitudes directes.
  \item Résoudre les exercices type examen national : nature de triangles, transformations planes, racines de l'unité.
\end{itemize}
\end{tcolorbox}

\vspace{4mm}

% ------------------- SECTION 1 -------------------
\section{Forme trigonométrique et exponentielle}

\begin{definition}[Argument]
Soit $z\in\mathbb{C}^*$, d'image $M$. Un \textbf{argument} de $z$, noté $\arg(z)$, est une mesure en radians de l'angle orienté $\bigl(\vec u\,;\overrightarrow{OM}\bigr)$. Il est défini \textbf{modulo $2\pi$} : si $\theta$ est un argument de $z$, tous les arguments de $z$ sont $\theta+2k\pi$, $k\in\mathbb{Z}$.
\end{definition}

\begin{propriete}[Forme trigonométrique]
Tout $z\in\mathbb{C}^*$ s'écrit
\[
  z = r(\cos\theta+i\sin\theta), \qquad r=|z|>0,\ \theta=\arg(z)\ [2\pi].
\]
Réciproquement, si $z=r(\cos\theta+i\sin\theta)$ avec $r>0$, alors $|z|=r$ et $\arg(z)=\theta\ [2\pi]$.
\end{propriete}

\begin{demo}
Si $z=a+ib$, $r=|z|=\sqrt{a^2+b^2}>0$. Posons $\theta$ tel que $\cos\theta=\frac ar$, $\sin\theta=\frac br$ (possible car $\left(\frac ar\right)^2+\left(\frac br\right)^2=\frac{a^2+b^2}{r^2}=1$, ce couple est bien sur le cercle trigonométrique). Alors $z=a+ib=r\cos\theta+ir\sin\theta=r(\cos\theta+i\sin\theta)$, et $\theta$ est bien l'angle $(\vec u;\overrightarrow{OM})$.
\end{demo}

\begin{definition}[Notation exponentielle]
Pour $\theta\in\mathbb{R}$, on note $e^{i\theta}=\cos\theta+i\sin\theta$ (notation d'Euler). Ainsi tout $z\in\mathbb{C}^*$ s'écrit sous \textbf{forme exponentielle} :
\[
  z = re^{i\theta}, \qquad r=|z|,\ \theta=\arg(z).
\]
\end{definition}

\begin{propriete}[Cohérence de la notation exponentielle]
\[
  e^{i\theta}\times e^{i\theta'} = e^{i(\theta+\theta')}, \qquad
  \frac{1}{e^{i\theta}} = e^{-i\theta}, \qquad
  \left(e^{i\theta}\right)^n = e^{in\theta}\ (n\in\mathbb{Z}), \qquad
  \overline{e^{i\theta}} = e^{-i\theta}, \qquad
  \bigl|e^{i\theta}\bigr| = 1.
\]
\end{propriete}

\begin{demo}
$e^{i\theta}e^{i\theta'} = (\cos\theta+i\sin\theta)(\cos\theta'+i\sin\theta')$. En développant et en utilisant les formules d'addition $\cos(\theta+\theta')=\cos\theta\cos\theta'-\sin\theta\sin\theta'$ et $\sin(\theta+\theta')=\sin\theta\cos\theta'+\cos\theta\sin\theta'$, on obtient exactement $\cos(\theta+\theta')+i\sin(\theta+\theta')=e^{i(\theta+\theta')}$.
\end{demo}

% ------------------- SECTION 2 -------------------
\section{Argument : propriétés et calculs}

\begin{propriete}[Propriétés de l'argument -- modulo $2\pi$]
Pour $z,z'\in\mathbb{C}^*$ :
\[
  \arg(zz') \equiv \arg(z)+\arg(z') \ [2\pi], \qquad
  \arg\left(\frac1z\right) \equiv -\arg(z) \ [2\pi],
\]
\[
  \arg\left(\frac{z}{z'}\right) \equiv \arg(z)-\arg(z') \ [2\pi], \qquad
  \arg(z^n) \equiv n\arg(z)\ [2\pi]\ (n\in\mathbb{Z}), \qquad
  \arg(\bar z) \equiv -\arg(z)\ [2\pi],
\]
\[
  \arg(-z) \equiv \arg(z)+\pi\ [2\pi].
\]
\end{propriete}

\begin{demo}
$\arg(zz')\equiv\arg z+\arg z'\ [2\pi]$ : en écrivant $z=re^{i\theta}$, $z'=r'e^{i\theta'}$, on a $zz'=rr'e^{i(\theta+\theta')}$, qui est bien la forme exponentielle de $zz'$ (module $rr'>0$, argument $\theta+\theta'$).
\end{demo}

\begin{propriete}[Formule de Moivre]
Pour tout $\theta\in\mathbb{R}$ et $n\in\mathbb{Z}$ :
\[
  (\cos\theta+i\sin\theta)^n = \cos(n\theta)+i\sin(n\theta).
\]
\end{propriete}

\begin{demo}
Conséquence immédiate de $\left(e^{i\theta}\right)^n=e^{in\theta}$.
\end{demo}

\begin{propriete}[Formules d'Euler]
Pour tout $\theta\in\mathbb{R}$ :
\[
  \cos\theta = \frac{e^{i\theta}+e^{-i\theta}}{2}, \qquad \sin\theta = \frac{e^{i\theta}-e^{-i\theta}}{2i}.
\]
\end{propriete}

\begin{demo}
On additionne (resp. soustrait) $e^{i\theta}=\cos\theta+i\sin\theta$ et $e^{-i\theta}=\cos\theta-i\sin\theta$, puis on isole $\cos\theta$ (resp. $\sin\theta$).
\end{demo}

\begin{exemple}[Linéarisation avec les formules d'Euler]
Linéarisons $\cos^3\theta$.
\[
  \cos^3\theta = \left(\frac{e^{i\theta}+e^{-i\theta}}{2}\right)^3
  = \frac{e^{3i\theta}+3e^{i\theta}+3e^{-i\theta}+e^{-3i\theta}}{8}
  = \frac{2\cos(3\theta)+6\cos\theta}{8} = \frac{\cos(3\theta)+3\cos\theta}{4}.
\]
\end{exemple}

\begin{exemple}[Calcul d'argument]
Déterminons module et argument de $z=1+i\sqrt3$.
\[
  |z| = \sqrt{1+3}=2, \qquad \cos\theta=\frac12,\ \sin\theta=\frac{\sqrt3}2 \Rightarrow \theta=\frac\pi3.
\]
\textbf{Conclusion} : $z=2e^{i\pi/3}$.
\end{exemple}

\begin{exemple}[Produit et quotient d'arguments]
Soient $z_1=2e^{i\pi/4}$, $z_2=3e^{-i\pi/6}$. Calculons $|z_1z_2|$, $\arg(z_1z_2)$, et $z_1/z_2$ sous forme exponentielle.
\[
  z_1z_2 = 6e^{i(\pi/4-\pi/6)} = 6e^{i\pi/12}, \qquad \frac{z_1}{z_2} = \frac23e^{i(\pi/4+\pi/6)} = \frac23e^{i5\pi/12}.
\]
\end{exemple}

% ------------------- SECTION 3 -------------------
\section{Racines $n$-ièmes}

\begin{propriete}[Racines $n$-ièmes de l'unité]
Pour $n\in\mathbb{N}^*$, les solutions de $z^n=1$ dans $\mathbb{C}$ sont les $n$ nombres
\[
  z_k = e^{i\frac{2k\pi}{n}}, \qquad k=0,1,\ldots,n-1.
\]
\end{propriete}

\begin{demo}
Cherchons $z=re^{i\theta}$ tel que $z^n=1$. Alors $r^ne^{in\theta}=1=1\cdot e^{i0}$, donc par unicité (module, argument modulo $2\pi$) : $r^n=1\Rightarrow r=1$ (car $r>0$), et $n\theta\equiv0\ [2\pi]\Rightarrow\theta\equiv0\ \left[\frac{2\pi}{n}\right]$, soit $\theta=\frac{2k\pi}{n}$ pour $k\in\mathbb{Z}$. En restreignant $k$ à $\{0,1,\ldots,n-1\}$, on obtient exactement $n$ solutions distinctes (les valeurs de $k$ hors de cet intervalle redonnent les mêmes points modulo $2\pi$).
\end{demo}

\begin{remarque}[Interprétation géométrique]
Les images des racines $n$-ièmes de l'unité forment un \textbf{polygone régulier à $n$ côtés} inscrit dans le cercle trigonométrique, avec un sommet en $1$.
\end{remarque}

\begin{propriete}[Racines $n$-ièmes d'un complexe quelconque]
Soit $Z=Re^{i\Phi}\in\mathbb{C}^*$ ($R>0$). Les solutions de $z^n=Z$ sont
\[
  z_k = R^{1/n}\,e^{i\left(\frac{\Phi}{n}+\frac{2k\pi}{n}\right)}, \qquad k=0,1,\ldots,n-1.
\]
\end{propriete}

\begin{exemple}[Racines cubiques de l'unité]
Résolvons $z^3=1$. Solutions : $z_0=1$, $z_1=e^{i2\pi/3}=-\frac12+i\frac{\sqrt3}2$, $z_2=e^{i4\pi/3}=-\frac12-i\frac{\sqrt3}2$.

\smallskip
On vérifie $1+z_1+z_2=0$ (propriété générale : la somme des racines $n$-ièmes de l'unité, pour $n\geq2$, est nulle).
\end{exemple}

\begin{exemple}[Racines quatrièmes d'un complexe]
Résolvons $z^4=-16$.

\smallskip
$-16=16e^{i\pi}$, donc $R=16$, $\Phi=\pi$. $R^{1/4}=2$.
\[
  z_k = 2\,e^{i\left(\frac\pi4+\frac{k\pi}2\right)}, \qquad k=0,1,2,3.
\]
Soit $z_0=2e^{i\pi/4}=\sqrt2+i\sqrt2$, $z_1=2e^{i3\pi/4}=-\sqrt2+i\sqrt2$, $z_2=-\sqrt2-i\sqrt2$, $z_3=\sqrt2-i\sqrt2$.
\end{exemple}

% ------------------- SECTION 4 -------------------
\section{Applications géométriques}

\begin{propriete}[Angle orienté et alignement]
Pour $A,B,C$ trois points d'affixes deux à deux distinctes $a,b,c$ :
\[
  \bigl(\overrightarrow{AB}\,;\overrightarrow{AC}\bigr) \equiv \arg\left(\frac{c-a}{b-a}\right) \ [2\pi].
\]
\begin{itemize}[leftmargin=6mm, itemsep=0.5mm]
  \item $A,B,C$ \textbf{alignés} $\iff \dfrac{c-a}{b-a}\in\mathbb{R}^*$.
  \item $(AB)\perp(AC) \iff \dfrac{c-a}{b-a}\in i\mathbb{R}^*$.
\end{itemize}
\end{propriete}

\begin{demo}
Notons $z_1=b-a$ (affixe de $\overrightarrow{AB}$) et $z_2=c-a$ (affixe de $\overrightarrow{AC}$). On a $\dfrac{z_2}{z_1} = \dfrac{|z_2|}{|z_1|}e^{i(\arg z_2-\arg z_1)}$, et $\arg z_2-\arg z_1 \equiv (\vec u;\overrightarrow{AC})-(\vec u;\overrightarrow{AB}) \equiv (\overrightarrow{AB};\overrightarrow{AC})\ [2\pi]$ (relation de Chasles sur les angles). D'où le résultat : le module de $\frac{c-a}{b-a}$ vaut $\frac{AC}{AB}$ et son argument est l'angle $(\overrightarrow{AB};\overrightarrow{AC})$.

\smallskip
$A,B,C$ alignés $\iff (\overrightarrow{AB};\overrightarrow{AC})\equiv0\ [\pi] \iff \arg\left(\frac{c-a}{b-a}\right)\equiv0\ [\pi] \iff \dfrac{c-a}{b-a}\in\mathbb{R}^*$.
\end{demo}

\begin{propriete}[Critère de cocyclicité -- admis]
Soient $A,B,C,D$ quatre points du plan, deux à deux distincts, d'affixes respectives $a,b,c,d$, avec $C,D\notin(AB)$. Les points $A,B,C,D$ sont \textbf{cocycliques} (situés sur un même cercle) \textbf{ou alignés} si et seulement si le \textbf{birapport}
\[
  \frac{(a-c)(b-d)}{(a-d)(b-c)}
\]
est un nombre \textbf{réel}.
\end{propriete}

\begin{remarque}[Justification (théorème de l'angle inscrit)]
$C$ et $D$ « voient » le segment $[AB]$ sous le même angle, ou sous des angles supplémentaires, c'est-à-dire $(\overrightarrow{CA};\overrightarrow{CB}) \equiv (\overrightarrow{DA};\overrightarrow{DB})\ [\pi]$, si et seulement si $A,B,C,D$ sont cocycliques (théorème de l'angle inscrit, cas dégénéré : alignement). Cette égalité d'angles se traduit, via l'argument d'un quotient de différences, par le caractère réel du birapport ci-dessus.
\end{remarque}

\begin{exemple}[Vérifier une cocyclicité]
Montrons que $A(0)$, $B(2)$, $C(1+i)$, $D(1-i)$ sont cocycliques.

\smallskip
Calculons le birapport :
\[
  \frac{(a-c)(b-d)}{(a-d)(b-c)} = \frac{\bigl(0-(1+i)\bigr)\bigl(2-(1-i)\bigr)}{\bigl(0-(1-i)\bigr)\bigl(2-(1+i)\bigr)} = \frac{(-1-i)(1+i)}{(-1+i)(1-i)}.
\]
Numérateur : $(-1-i)(1+i)=-1-i-i-i^2=-2i$. Dénominateur : $(-1+i)(1-i)=-1+i+i-i^2=2i$.
\[
  \frac{-2i}{2i} = -1 \in \mathbb{R}.
\]
\textbf{Conclusion} : le birapport est réel, donc $A,B,C,D$ sont cocycliques ou alignés ; comme $C,D\notin(AB)$ (l'axe réel), ils sont bien \textbf{cocycliques} (on vérifie d'ailleurs que les quatre points sont à distance $1$ du point d'affixe $1$ : ils sont sur le cercle de centre $\Omega(1)$ et de rayon $1$).
\end{exemple}

\begin{propriete}[Expressions complexes des transformations usuelles]
Soit $M(z)$ un point du plan, d'image $M'(z')$ par une transformation. Dans un repère orthonormé direct :
\begin{itemize}[leftmargin=6mm, itemsep=0.5mm]
  \item \textbf{Translation} de vecteur $\vec w$ d'affixe $b$ : \quad $z'=z+b$.
  \item \textbf{Homothétie} de centre $\Omega(\omega)$ et de rapport $k\in\mathbb{R}^*$ : \quad $z'-\omega=k(z-\omega)$, soit $z'=kz+(1-k)\omega$.
  \item \textbf{Rotation} de centre $\Omega(\omega)$ et d'angle $\theta$ : \quad $z'-\omega=e^{i\theta}(z-\omega)$, soit $z'=e^{i\theta}z+\bigl(1-e^{i\theta}\bigr)\omega$.
\end{itemize}
\end{propriete}

\begin{demo}[Rotation]
Par définition géométrique, une rotation de centre $\Omega$ et d'angle $\theta$ envoie $M$ sur $M'$ tel que $\Omega M'=\Omega M$ et $(\overrightarrow{\Omega M}\,;\overrightarrow{\Omega M'})\equiv\theta\,[2\pi]$. En termes d'affixes (grâce à la propriété de l'angle orienté vue plus haut, appliquée à $\overrightarrow{\Omega M}$ et $\overrightarrow{\Omega M'}$) :
\[
  \frac{z'-\omega}{z-\omega} = \left|\frac{z'-\omega}{z-\omega}\right|e^{i\theta} = 1\times e^{i\theta} = e^{i\theta}
\]
(le module vaut $1$ car $\Omega M'=\Omega M$), d'où $z'-\omega=e^{i\theta}(z-\omega)$. Le cas de l'homothétie se traite de façon analogue, avec un rapport réel $k$ au lieu de $e^{i\theta}$ (angle nul si $k>0$, angle $\pi$ si $k<0$).
\end{demo}

\begin{exemple}[Reconnaître une transformation à partir de son expression complexe]
Identifions la transformation $z'=iz+2$.

\smallskip
On a $a=i$, avec $|a|=1$ et $\arg(a)=\frac\pi2$ : c'est une \textbf{rotation} d'angle $\frac\pi2$. Centre : point fixe $\omega=\dfrac{b}{1-a}=\dfrac{2}{1-i}=\dfrac{2(1+i)}{(1-i)(1+i)}=\dfrac{2(1+i)}{2}=1+i$.

\smallskip
\textbf{Conclusion} : rotation de centre $\Omega(1+i)$ et d'angle $\dfrac\pi2$.
\end{exemple}

\begin{propriete}[Similitude directe]
Une \textbf{similitude directe} de rapport $k>0$, d'angle $\theta$, de centre $\Omega$ (affixe $\omega$) est l'application qui à $M(z)$ associe $M'(z')$ tel que
\[
  z' - \omega = ke^{i\theta}(z-\omega).
\]
Réciproquement, toute transformation $z'=az+b$ avec $a\in\mathbb{C}^*$ est une similitude directe de rapport $|a|$, d'angle $\arg(a)$, dont le centre (si $a\ne1$) est le point fixe $\omega=\dfrac{b}{1-a}$.
\end{propriete}

\begin{exemple}[Nature d'un triangle via le rapport complexe]
Soient $A(0)$, $B(1+i)$, $C(2i)$. Déterminons la nature du triangle $ABC$.
\[
  \frac{c-a}{b-a} = \frac{2i}{1+i} = \frac{2i(1-i)}{(1+i)(1-i)} = \frac{2i+2}{2} = 1+i.
\]
$|1+i|=\sqrt2$ et $\arg(1+i)=\frac\pi4$. Donc $\dfrac{AC}{AB}=\sqrt2$ et $(\overrightarrow{AB};\overrightarrow{AC})=\frac\pi4$.

\smallskip
\textbf{Conclusion} : le triangle $ABC$ n'est ni équilatéral ni isocèle en $A$ directement lisible, mais l'angle en $A$ vaut $45°$ ; pour un triangle \textbf{rectangle isocèle} en $A$ on aurait attendu $\left|\frac{c-a}{b-a}\right|=1$ et argument $\pm\frac\pi2$ -- ici ce n'est pas le cas, donc ce triangle n'est ni rectangle ni isocèle en $A$.
\end{exemple}

% ------------------- SECTION 5 -------------------
\section{Exercices corrigés -- niveau examen national SM}

\begin{exemple}[Exercice 1 -- forme exponentielle et puissance]
Écrire $z=\left(\sqrt3-i\right)^6$ sous forme algébrique, en passant par la forme exponentielle.

\smallskip
\textbf{Corrigé.} $|\sqrt3-i|=\sqrt{3+1}=2$ ; $\cos\theta=\frac{\sqrt3}2$, $\sin\theta=-\frac12\Rightarrow\theta=-\frac\pi6$. Donc $\sqrt3-i=2e^{-i\pi/6}$.
\[
  z = 2^6e^{-i6\pi/6} = 64e^{-i\pi} = 64\bigl(\cos(-\pi)+i\sin(-\pi)\bigr) = 64\times(-1) = -64.
\]
\end{exemple}

\begin{exemple}[Exercice 2 -- triangle équilatéral direct]
Montrer que le triangle $ABC$ d'affixes $a=0$, $b=1$, $c=e^{i\pi/3}$ est équilatéral direct.

\smallskip
\textbf{Corrigé.} $\dfrac{c-a}{b-a} = \dfrac{e^{i\pi/3}}{1} = e^{i\pi/3}$, de module $1$ et d'argument $\frac\pi3$.

\smallskip
Module $1$ $\Rightarrow AC=AB$ : triangle isocèle en $A$. Argument $\frac\pi3$ $\Rightarrow$ angle en $A$ vaut $60°$. Un triangle isocèle avec un angle au sommet de $60°$ est \textbf{équilatéral} (et l'orientation directe est donnée par l'argument positif).
\end{exemple}

\begin{exemple}[Exercice 3 -- équation avec racines $n$-ièmes et somme]
Résoudre $z^5=32i$, puis calculer la somme des cinq solutions.

\smallskip
\textbf{Corrigé.} $32i=32e^{i\pi/2}$. $R^{1/5}=32^{1/5}=2$.
\[
  z_k = 2e^{i\left(\frac{\pi}{10}+\frac{2k\pi}5\right)}, \qquad k=0,1,2,3,4.
\]
\emph{Somme} : les $z_k$ sont les racines cinquièmes de $32i$, obtenues en multipliant les racines cinquièmes de l'unité par $2e^{i\pi/10}$ ; comme la somme des racines cinquièmes de l'unité est nulle (pour $n\geq2$), la somme des $z_k$ est également nulle : $\displaystyle\sum_{k=0}^4z_k=0$.
\end{exemple}

\begin{exemple}[Exercice 4 -- similitude directe]
Soit $f$ la transformation d'expression complexe $z'=(1+i)z-2i$. Déterminer la nature et les éléments caractéristiques de $f$.

\smallskip
\textbf{Corrigé.} $a=1+i$, $b=-2i$. $|a|=\sqrt2$, $\arg(a)=\frac\pi4$ : $f$ est une \textbf{similitude directe} de rapport $\sqrt2$ et d'angle $\frac\pi4$.

\smallskip
Centre : point fixe $\omega=\dfrac{b}{1-a}=\dfrac{-2i}{1-(1+i)}=\dfrac{-2i}{-i}=2$.

\smallskip
\textbf{Conclusion} : $f$ est la similitude directe de centre $\Omega(2)$, de rapport $\sqrt2$, d'angle $\frac\pi4$.
\end{exemple}

\begin{exemple}[Exercice 5 -- linéarisation et intégrale (lien inter-chapitres)]
Utiliser la linéarisation de $\cos^2\theta$ (via Euler) pour retrouver $\displaystyle\int_0^{\pi}\cos^2\theta\,d\theta$.

\smallskip
\textbf{Corrigé.} $\cos^2\theta=\left(\dfrac{e^{i\theta}+e^{-i\theta}}2\right)^2=\dfrac{e^{2i\theta}+2+e^{-2i\theta}}4=\dfrac{2\cos(2\theta)+2}4=\dfrac{1+\cos(2\theta)}2$.
\[
  \int_0^\pi\cos^2\theta\,d\theta = \int_0^\pi\frac{1+\cos(2\theta)}2\,d\theta = \left[\frac\theta2+\frac{\sin(2\theta)}4\right]_0^\pi = \frac\pi2.
\]
\end{exemple}

% ------------------- SECTION 6 -------------------
\section{Méthodes et pièges : la boîte à outils de l'examen}

\subsection{Ce qu'on te demande, ce que tu fais}

\begin{center}
\renewcommand{\arraystretch}{1.45}
\sffamily\small
\begin{tabular}{>{\raggedright\arraybackslash}p{7.2cm} >{\raggedright\arraybackslash}p{8cm}}
\rowcolor{qtealdark} \color{white}\bfseries Ce qu'on te demande & \color{white}\bfseries Ton réflexe \\
\rowcolor{tealpale} Trouver module et argument & $r=|z|=\sqrt{a^2+b^2}$, puis $\cos\theta=\frac ar$, $\sin\theta=\frac br$. \\
Puissance $n$-ième & Forme exponentielle $z=re^{i\theta}$, puis $z^n=r^ne^{in\theta}$ (Moivre). \\
\rowcolor{tealpale} Linéariser $\cos^n\theta$ ou $\sin^n\theta$ & Formules d'Euler, développer le binôme, regrouper les termes conjugués. \\
Racines $n$-ièmes de $Z=Re^{i\Phi}$ & $z_k=R^{1/n}e^{i(\Phi/n+2k\pi/n)}$, $k=0,\ldots,n-1$. \\
\rowcolor{tealpale} Alignement / orthogonalité de points & Étudier $\dfrac{c-a}{b-a}$ : réel (aligné) ou imaginaire pur (orthogonal). \\
Montrer que 4 points sont cocycliques & Calculer le birapport $\dfrac{(a-c)(b-d)}{(a-d)(b-c)}$ : réel $\Rightarrow$ cocycliques ou alignés. \\
Nature d'un triangle & Module de $\frac{c-a}{b-a}$ = rapport des côtés, argument = angle en $A$. \\
\rowcolor{tealpale} Identifier une similitude $z'=az+b$ & Rapport $|a|$, angle $\arg(a)$, centre $\omega=\frac{b}{1-a}$ (si $a\ne1$). \\
Reconnaître translation/rotation/homothétie & $|a|=1,\ a\ne1$ : rotation d'angle $\arg a$ ; $a\in\mathbb{R}^*\setminus\{1\}$ : homothétie de rapport $a$ ; $a=1$ : translation. \\
\end{tabular}
\end{center}

\subsection{Les pièges qui coûtent des points}

\begin{remarque}[Erreurs relevées chaque année par les correcteurs]
\begin{itemize}[leftmargin=6mm, itemsep=0.5mm]
  \item Oublier « modulo $2\pi$ » en manipulant les arguments : $\arg(zz')\equiv\arg z+\arg z'\,[2\pi]$, pas une égalité stricte.
  \item Confondre $\left(\frac{c-a}{b-a}\right)$ (angle en $A$) avec un autre rapport (bien identifier le sommet).
  \item Oublier de vérifier $a,b,c$ deux à deux \textbf{distincts} avant d'utiliser les formules d'angle orienté.
  \item Se tromper dans les racines $n$-ièmes en ne balayant pas $k=0$ à $n-1$ (oubli d'une racine ou doublon).
  \item Écrire $e^{i\theta}\times e^{i\theta'}=e^{i\theta\theta'}$ au lieu de $e^{i(\theta+\theta')}$.
  \item Pour la similitude, oublier que le centre n'est défini que si $a\ne1$ (sinon $f$ est une translation, pas une similitude au sens strict).
\end{itemize}
\end{remarque}

% ------------------- RESUME -------------------
\vspace{4mm}
\begin{tcolorbox}[
  enhanced, boxrule=0pt, arc=2mm,
  interior style={left color=qtealdark, right color=qteal},
  left=6mm, right=6mm, top=4mm, bottom=4mm,
  fontupper=\color{white}\sffamily,
  title={\sffamily\bfseries\color{white}\ding{72}~~À retenir},
  colbacktitle=qtealdark, coltitle=white, boxrule=0pt
]
\begin{itemize}[leftmargin=6mm, label={\color{qamber}\ding{212}}, itemsep=1mm]
  \item $z=re^{i\theta}$, $r=|z|$, $\theta=\arg z\,[2\pi]$ ; $\arg(zz')\equiv\arg z+\arg z'$ ; Moivre : $(\cos\theta+i\sin\theta)^n=\cos n\theta+i\sin n\theta$.
  \item Euler : $\cos\theta=\frac{e^{i\theta}+e^{-i\theta}}2$, $\sin\theta=\frac{e^{i\theta}-e^{-i\theta}}{2i}$.
  \item Racines $n$-ièmes de $Re^{i\Phi}$ : $z_k=R^{1/n}e^{i(\Phi/n+2k\pi/n)}$ ; racines de l'unité : polygone régulier, somme nulle.
  \item Angle orienté $(\overrightarrow{AB};\overrightarrow{AC})=\arg\left(\frac{c-a}{b-a}\right)$ ; alignement $\iff$ rapport réel ; orthogonalité $\iff$ rapport imaginaire pur.
  \item Cocyclicité de $A,B,C,D$ : birapport $\dfrac{(a-c)(b-d)}{(a-d)(b-c)}\in\mathbb{R}$ (ou alignement, cas dégénéré).
  \item Translation : $z'=z+b$ ; homothétie $(\Omega,k)$ : $z'-\omega=k(z-\omega)$ ; rotation $(\Omega,\theta)$ : $z'-\omega=e^{i\theta}(z-\omega)$.
  \item Similitude directe $z'=az+b$ : rapport $|a|$, angle $\arg a$, centre $\frac{b}{1-a}$ ; cas particuliers : $a=1$ translation, $a\in\mathbb{R}^*$ homothétie, $|a|=1$ rotation.
\end{itemize}
\end{tcolorbox}

\end{document}
